高效准确的锂离子电池健康状态估计，对保障电池安全运行和支持资源受限电池管理系统的在线应用至关重要。然而，现有估计方法往往依赖跨电池稳定性有限的健康指标和资源开销较大的模型结构。为此，本文提出一种轻量化的SOH估计框架。首先，通过组级健康指标选择算法，从多源候选指标中选取跨电池稳健的健康指标，并减少特征冗余。随后，构建轻量化预测网络MS-AgentNet。该网络主要集成局部—全局融合注意力模块，通过小核深度可分离卷积与ReLU²智能体注意力协同捕获局部退化特征与长程退化依赖，同时将传统Transformer的二次注意力复杂度降至线性。此外，大核深度可分离卷积用于提取长尺度退化特征，并与小核卷积共同以较低参数开销融合多尺度和多通道特征，增强特征多样性。我们在四个涵盖不同化学体系和运行条件的公开电池数据集上，将所提模型与多种主流深度学习模型进行了比较。实验结果表明，MS-AgentNet总体上取得了更优的综合性能，同时保持了较低的计算与存储开销，进一步体现了其面向资源受限电池管理系统的轻量化优势。

\textbf{关键词：} 锂离子电池；健康状态估计；线性复杂度；深度特征融合；轻量化深度学习
